% ================================================================================================ %
% Hilfe zum Einstieg in LaTeX
% ================================================================================================ %

% Hier sind ein paar nützliche Links um LaTeX besser lernen zu können und ein paar die es sich zu merken einfach lohnt.

% Ein guter Einstieg ist die Overleaf Dokumentation, diese nimmt dich gut an die Hand.
    % https://de.overleaf.com/learn

% Eine weitere gute Quelle ist "The LaTeX Project", diese bieten einen weiteren gute Dokumentation (auch in De) um hier einen umfassenden Einstig zu finden.
    % https://www.latex-project.org/help/documentation/

% Wenn man schnell LaTeX lernen will empfehle ich die 5 Artikel von Overleaf zum Thesis schreiben
    % https://de.overleaf.com/learn/latex/How_to_Write_a_Thesis_in_LaTeX_(Part_1)%3A_Basic_Structure

% oder die 18 kurze und anschauliche Kapitel von Latex-tutorials.com oder iht Quikstart-Guide
    % https://latex-tutorial.com/tutorials/
    % https://latex-tutorial.com/quick-start/
    
% Auch sinnvoll ist es bei Problemen zu Googeln oder ne KI-Chatbox zu fragen, die meisten können mit LaTeX umgehen und geben aus meiner Erfahrung relativ Sinnvolle antworten

% ======= Resourcen, die Sinnvoll sind %
% Hyperref Anpassungen
    % https://www.namsu.de/Extra/pakete/Hyperref.html
% Biblatex Quide zum einstieg in die Quellen
    % https://de.overleaf.com/learn/latex/Bibliography_management_with_biblatex
% nutzung von Farben in LaTeX (mit xColor)
    % https://latex-tutorial.com/color-latex/
    


\documentclass[             % Basiseinstellungen des Dokuments
12pt,                           % Schriftgröße
a4paper                         % Blattformat
]{report}                       % Art des Textes (Report, Book, Article, ...)
\usepackage[utf8]{inputenc} % Unicode Standart für Sonderzeichen
\usepackage[T1]{fontenc}    % Hinzufügen von mehr Seitenumbrüchen
\usepackage[]{fancyhdr}     % Hinzufügen von mehr Seitenstyles
\pagestyle{fancy}           % Änderung zum "fancy" Seitenstyles
\usepackage[ngerman]{babel} % Deutsches Sprachpaket mit neue Rechtschreibung, d.\,h. (Silbentrennung)
\usepackage[dvipsnames, x11names ]{xcolor} % zusätzliche Farben
\usepackage{multirow}       % Hinzufügen von Tabellen
\usepackage{array}          % Mehr optionen bei Tabellen 
\usepackage{graphicx}       % besseres Hinzufügen und Formatieren von Bildern 
\usepackage[iso, german]{isodate}      % formatieren von Datum, Zeit und Zeitzonen
\usepackage{hyperref}       % Hinzufügen von internen und externen Links
\usepackage{lipsum}         % fügt Text hinzu (Lorem Ipsum)       
  \usepackage[              % fügt Biblatex für vernünftige Quellen ein
    backend=biber,              % Formatierung vom Backend der Quellendokumente
    style=numeric,              % Definiert die Formatierung 
  ]{biblatex}
\usepackage{colortbl}       % Ermöglicht es Tabellen Farbig zu gestalten
\usepackage{movie15}        % Ermöglicht das einbinden von .gif datein durch den befehl
\usepackage{awesomebox}
\usepackage{makecell}       % Ermöglicht linebreaks in Tabellen durch \makecell{} und Überschriften dank \thead{} 

% ================================================================================================ %
% Anpassungen & Befehle
% ================================================================================================ %
% Eine Anmerkung zum Vertikalen Platz


\hypersetup{                % Einstellung von Hyperref
% ======= Dokumentoptionen %
pdftoolbar=true,	% Anzeigen der Acrobat toolbar oder nicht
pdfmenubar=true,	% Anzeigen des Acrobat menu oder nicht
pdftitle={Zusammenfassung Semester 1 Fakultät DM},	% Titel
pdfsubject={lernen},	% Um was geht es
pdfauthor={Nils J. Hack, },	% Autor bzw. Autoren
pdfkeywords={Hochschule Furtwangen, Fakultät Digitale Medien, SoSe24},	% Stichwort1, Stichwort2 ...
pdfcreator={Overleaf},	% Mit welcher Anwendung i.d.R. pdflatex
pdfproducer={Text},	% LaTeX with hyperref
% ======= %
bookmarks=true,	            % erstellt Bookmarks
bookmarksopen=true,        % Anzeigen der Bookmarks beim Öffen des Dokuments
bookmarksnumbered=false,	% Anschnittsnummer anzeigen
% ======= %
    colorlinks=true,            % Farblinks an
    linkcolor=YellowGreen,      % Farbe von internen Verweise
    citecolor=Cerulean,         % Farbe von Zitaten
    urlcolor=RedViolet,         % Farbe zu externen Links
    }

% ============= %
\renewcommand{\footrulewidth}{0.4pt}    % 
% ============= %
\setlength\parindent{0pt}
% ============= %
\setcounter{secnumdepth}{5}
% ============= %

% ======= Quellenbibliotheken hinzufügen %
\addbibresource{Quellen/Medientechnik.bib}


% ======= Formatierungen des Seitenlayouts %
\fancyhead[R]{\rightmark}   % Kopfzeile mit aktueller Section rechts
\fancyhead[L]{\leftmark}    % Kopfzeile mit aktuellem Kapitel oben links
\fancyfoot[R]{\hyperref[TOC]{\textcolor{YellowGreen}{Inhaltsverzeichnis}}}  % fügt unten rechts einen link zum Inhaltsverzeichnis hinzu
\fancyfoot[L]{SoSe24 DM Sem1} % fügt unten links den Text hinzu


% ================================================================================================ %
% Dokument beginnt
% ================================================================================================ %

\begin{document}

% ======= Titelseite anpassen %
\vspace{-5cm}
\title{\vspace{-5cm}\begin{center}
    \includegraphics[height=100px]{Bilder/HFU_logo.png}
\end{center}\vspace{3cm}Zusammenfassung\\ Semester 1\\ Fakultät DM}
\author{Zusammengefasst von:\\\emph{Nils Hack \& Daniel Georg}}
\date{Im Zeitraum:\\12. Juli 2024 - \today}

% ======= Titelseite %
\maketitle
\setcounter{page}{2} % Setzt Seitenzahl auf 2 weil Deckblatt geskippt wird
\newpage

% ======= Ínhaltsverzeichnis %
\hypersetup{linkcolor=black}    % Setzt die Farbe des Inhaltsverzeichnis zu schwarz, überschreibt das Seitenlayout
\tableofcontents\label{TOC}     % fügt ein label ein um dieses zu Verlinken
\newpage

% ======= Inhalte %             % Durch das hinzufügen von einem "%" die ungewünchten Bereiche ausklammern
\section{Vorwort}
    \subsection{Formattierung}
Unten Rechts am Rand unter \textcolor{YellowGreen}{Inhaltsverzeichnis} ist ein Link der dich wieder nach oben zum Inhaltsverzeichnis bringt.\\~\\
Zudem wird obenlinks das Kapitel und oben rechts das unterkapitel angezeigt um die Orientierung zu verbessern.\\~\\
Im folgenem Dokument werden \\
\textcolor{YellowGreen}{interne Link} in YellowGreen,   \\   
\textcolor{Cerulean}{Quellenverlinkungen} in SeaGreen und \\         
\textcolor{RedViolet}{externe Links} in  RedViolet dargesellt.
    \subsection{Zum Inhalt}
Dies hier zuammengestelten Seiten versuchen ein Gesamtbild aller Vorlesungen abzubilden, die im 1. Semester in der noch so genannten Fakultät Digitale Medien beigebrachten Inhalte näher zu bringen. Dies wird wahrscheinlich nicht im vollem Umfang möglich sein und ich bitte dies zu berücksichtigen. Auch werden wir Versuchen hier aktuelle Klausuraufgaben möglichst gut wiederzugeben, jedoch halten wir uns hier an dem gesetzlichen Rahmen und werden diese entsprechend nicht 1:1 veröffentlichen.\\~\\
Im inhaltlichen Fangen wir bei der SPO an, um einen Rahmen zu bieten und um darauf inhaltlich Aufbauen zu können. Wir Orientieren uns nach der Bennenung nach der SPO und entsprechend sind die Kapitel getaltet. Am Ende von jedem Kurs wollen wir noch kurz mögliche Klausuraufgaben erklären und wo es Sinn ergibt auch das Inhaltliche etwas erweitern und evtl. praktische Tipps näher zu bringen.\\~\\
\emph{Nils}
\chapter{AWBL Maier}

\section{SPO}

    Nachdem Studierende das Modul erfolgreich abgeschlossen haben, können sie:
    \begin{itemize}
        \item Wissen / Kenntnisse\\
            Das Mikro- und Makro-Umfeld von Medienbetrieben, und wie sie diese beeinflussen, benennen sowie den Stellenwert der Medienbranche in der Volkswirtschaft und Gesellschaft skizzieren.
            \newline
            Erklären, wie Medienunternehmen aus betriebswirtschaftlicher Sicht grundlegend funktionieren sowie die relevanten regulatorischen Bedingungen für das Medienmanagement kennen.
        \item Verstehen\\
            Verstehen, wie sich Medienbetrieb unterschiedlicher Art finanzieren sowie verstehen, welche Rechtsformen Medienbetriebe haben können.
            \newline
            Verstehen, welche strategischen und operativen Entscheidungen Medienunternehmen treffen müssen sowie welche Managementinstrumente Medienunternehmen benutzen (können).
        \item Anwenden\\
            Darlegen, in welchem volkswirtschaftlichen, politischen und regulatorischen Bezugsrahmen Medienbetriebe agieren. 
            \newline
            Benennen, wie einzelne Medienbetriebe ihren Markt bzw. ihre Branche definieren und wie sie dies in ihren Aktivitäten beeinflusst.
        \newpage
        \item Analyse\\
            Analysieren, wie Angebot und Nachfrage von Mediengütern zusammenspielen und wie dies von Medienbetrieben koordiniert wird sowie Investitionsentscheidungen in Medienbetrieben analysieren.
            \newline
            Analysieren, wie Medienunternehmen organisiert sind sowie welche Auswirkungen regulatorische Bedingungen auf Entscheidungen im Medienmanagement haben.
        \item Synthesis\\
            Allgemeine personalpolitische Maßnahmen auf Medienbetriebe übertragen.
        \item Evaluation\\
            Steuerliche Konsequenzen medienbetrieblicher Entscheidungen grob bewerten.
            \newline
            Medienbetriebliche Entscheidungen aus Sicht des Controllings bewerten.
    \end{itemize}
    
\newpage
\section{Allgemeine BWL}

\notebox{Aufgabe der Betriebswirtschaftslehre ist es, alles wirtschaftliche Handeln, das sich im Betrieb vollzieht, zu beschreiben und zu erklären und schließlich auf Grund der erkannten Regelmäßigkeiten und Gesetzmäßigkeiten des Betriebsprozesses wirtschaftliche Verfahren zur Realisierung praktischer betrieblicher Zielsetzungen zu entwickeln.}


\newpage\section{Unternehmensgründung}
    \subsection{Grundlagen}

\begin{itemize}
    \item Aus welchen Motiven Gründen Menschen Unternehmen?Was zeichnet einen guten Entrepeneur aus?
    \begin{itemize}
        \item Unzufriedenheit mit Unternehmen
        \item Eigenständigkeit / kein Chef
        \item mehr Freiheiten als Angestellte
        \item höheres Gehalt
        \item Sinnhaftigkeit
        \item eigene Idee Umsetzten
        \item Ausnutzung Marktlücke / Nische
        \item  geringe Auslastung in einem Bereich
    \end{itemize}
    \item Was zeichnet einen guten Entrepreneur aus?
    \begin{itemize}
        \item Kritikfähig
        \item Kompetenz
        \item Work-a-holic
        \item Auslagerung von Kompetenzen
        \item hartnäckig
        \item ausdauernd
        \item Einfühlungsvermögen
        \item Begeisterungsfähigkeit
        \item Überzeugungskraft
        \item Im Interesse des Unternehmens handeln
        \item Verbesserungsfähig
        \item Fehler eingestehen
        \item kein Arschloch
        \item fair sein / beurteilen
        \item nicht nachtragend
    \end{itemize}\newpage
    \item Was sind die Vor- und Nachteile einer Unternehmensgründung?
    \begin{itemize}
        \item höheres Risiko
        \item Insolvenzrisiko
        \item Zeitintensiv
        \item hoher Verwaltungsaufwand
        \item mehr Freiheiten
        \item möglicherweise keine Kunden
        \item höheres persönliches Risiko
        \item anfänglich geringe Gewinne (wenn überhaupt) 
    \end{itemize}
    \item Schritte einer Gründung
    \begin{itemize}
        \item Kapital
        \item Idee
        \item Investment falls nötig
        \item Angestellte
        \item Prototyp
        \item Rechtsform
        \item Business Plans
        \item Kundengewinnung
    \end{itemize}
    \item Faktoren
    \begin{itemize}
        \item Schritte
        \item siehe Risiken
    \end{itemize}
    \item Gründe zum Scheitern
    \begin{itemize}
        \item Schlechte Planung
        \item Unerwartete Ereignisse
        \item Verkalkulation
        \item zu hohe Steuervorzahlungen
        \item schlechte Markteinschätzung
        \item zu hohe Kosten
        \item Mentale Probleme
        \item zu schnelles Wachstum
        \item falsche Prognosen
        \item externe Einflüsse (Katastrophen, Kriege, ...) 
    \end{itemize}
\end{itemize}
    \subsection{Unternehmensplan}
    \subsection{Die 9 Schritte zur Selbstständigkeit}
        \subsubsection{Entscheidung für die Selbstständigkeit}
            \paragraph*{Motive für die Existenzgründung}
            \begin{itemize}
                \item Innovation
                \item Anerkennung
                \item Rollenverhalten
                \item Selbstverwirklichung
                \item Unabhängigkeit
                \item Wirtschaftlicher Erfolg
            \end{itemize}

            
            \paragraph*{Eigenschaften eines Gründers}
            \begin{itemize}
                \item Flexibilität
                \item Machbarkeitsüberzeugung
                \item Risikofreudigkeit
                \item Soziale Kompetenz
                \item Entschlussfreudigkeit
                \item Problemorientierung
                \item Wachstumsorientierung
                \item Durchhaltevermögen
                \item Unabhängigkeitsstreben
                \item Leistungsmotiv
            \end{itemize}
            \paragraph*{ Auslöser der Gründungsaktivität „Theory of planned behavior“}
                „Unternehmensgründungen sind kein spontanes Ereignis zu einem zufälligen Zeitpunkt, sondern das Ergebnis von situativen und kulturellen Faktoren.“\newpage
                \begin{center}
                Äußere oder innere Lebensumstände\\+\\
                Positive Bewertung der Selbstständigkeit\\+\\
                Persönliche hohe Handlungsbereitschaft\\=\\
                Wahrscheinlichkeit der Unternehmensgründung \end{center}
        \subsubsection{Zusammenstellung eines Teams}
            \paragraph{Gründe für eine Gründung im Team}~\\
            \begin{itemize}
                \item Ausgleich der vorhandenen Schwächen (Persönlichkeit, Kompetenz, Know-how)
                \item Größere Finanzkraft
                \item Geteiltes finanzielles Risiko 
                \item Gegenseitige Sparringspartner bei Ideen- und Entscheidungsfindung
            \end{itemize}
            \paragraph{Aber...}~\\
            \begin{itemize}
                \item Verlust an Autonomie
                \item Aufteilen der Erlöse auf mehrere Gründer anteilig an Beteiligung
            \end{itemize}
            \notebox{Eine durchdachte Zusammenstellung des Gründerteams ist entscheidend!}
            \paragraph{Merkmale einer guten Teamzusammenstellung}~\\
            \begin{itemize}
                \item Teile eine gemeinsamen Vision
                \item Gemeinsame Motivation
                \item Hohe Teamfähigkeit und gegenseitige Unterstützung
                \item Offene und regelmäßige Kommunikation
                \item Komplmentäre Eigenschaften und stärken
                \item Klare Vereinbarungen über Eigenstumsverhältnisse
                \item Klare Vereinbarungen über Rechte und Pflichten
                \item Klare Aufteilung der Zuständigkeiten
            \end{itemize}

            \paragraph{Eine häufige Aufteilung in einem Start-Up ist}
            \begin{itemize}
                \item Einer übernimmt die Entwicklung/Technische Leistung
                \item der andere übernimmt den Vertrieb / die kaufmännische Leitung
            \end{itemize}
        \subsubsection{Geschäftsidee entwickeln}
        \subsubsection{Geschäftsmodell konzipieren}
        \subsubsection{Businessplan aufstellen}
        \subsubsection{Finanzierung}
        \subsubsection{Unternehmen gründen}
        \subsubsection{Angebot vermarkten}
        \subsubsection{Erfolg hinterfragen}

    
\input{AWBL/2_AWBL_Prüfungsaufgaben}
\chapter{Medientechnik}

\section{SPO}
    \subsection{Audiotechnik}
        \subsubsection{Lernergebnisse:}
        Nachdem Studierende das Modul erfolgreich abgeschlossen haben, können sie:

        \begin{itemize}
            \item Wissen / Kenntnisse\\
                die AV-technischen Voraussetzungen der computerbasierten Medienproduktion kennen und beherrschen
            \item Verstehen\\
                die physikalischen AV-Grundlagen in computerbasierten Medienanwendungen in Beziehung setzen
            \item Anwenden\\
                die erworbenen theoretischen und technischen Kenntnisse auf konkrete Medienanwendungen übertragen
            \item Analyse\\
                Aufgabenstellungen in computerbasierten Medienproduktionen erkennen und analysieren sowie deren Durchführung planen
            \item Synthesis\\
                infache AV-Produktionen zusammen mit computergenerierten Inhalten durchführen.
            \item Evaluation\\
            etwaige Fehler im computerbasierten AV-Produktionsprozess erkennen und korrigieren.\\
            sicher mit AV-Produktionsequipment umgehen
        \end{itemize}
\newpage
\section{Audiotechnik Reusch}

\input{Medientechnik/Audiotechnik_Reusch/Einstieg_Reusch}
\input{Medientechnik/Audiotechnik_Reusch/Schallwellen_Reusch}
\input{Medientechnik/Audiotechnik_Reusch/Ohr_Reusch}
\input{Medientechnik/Audiotechnik_Reusch/Mikrofone_Reusch}
\input{Medientechnik/Audiotechnik_Reusch/Lokalisation_Reusch}
\input{Medientechnik/Audiotechnik_Reusch/Stereo_reusch}
\input{Medientechnik/Audiotechnik_Reusch/Mikrofonierung_Reusch}
\input{Medientechnik/Audiotechnik_Reusch/Anschlusstechnik_Reusch}
\input{Medientechnik/Audiotechnik_Reusch/DAW_Reusch}
\input{Medientechnik/Audiotechnik_Reusch/Aufzeichnung_Reusch}
\input{Medientechnik/Audiotechnik_Reusch/Praxisbeispiele}







\section{Videotechnik Hottong}

\input{Medientechnik/Videotechnik_Hottong/Licht_Hottong}
\input{Medientechnik/Videotechnik_Hottong/Objektive_Hottong}
\input{Medientechnik/Videotechnik_Hottong/Wahrnehmung_hottong}
\input{Medientechnik/Videotechnik_Hottong/Techn_Qualitaet_Hottong}
\input{Medientechnik/Videotechnik_Hottong/Farbmodelle_Hottong}
\input{Medientechnik/Videotechnik_Hottong/Generierung_Bilder_Hottong}
\input{Medientechnik/Videotechnik_Hottong/Bildaufnahmetechniken_Hottong}

\chapter{Gestaltung}
    \section{SPO}
    Nachdem Studierende das Modul erfolgreich abgeschlossen haben, können sie:
    \begin{itemize}
        \item Wissen / Kenntnisse\\
            Die Grundlagen gestalterischer Fragestellungen beurteilen.
            \newline
            Theorien zur Medienrezeption benennen.
        \item Verstehen\\
            Kreative Prozesse verstehen und selber erste Gestaltarbeiten anfertigen.
            \newline
            Verstehen, wo wir als Rezipienten und als Produzenten auf wissenschaftliche Erkenntnisse aufbauen können.
        \item Anwenden\\
            Erste Konzeptionen entwickeln und mit den Augen eines Gestalters Kreativarbeit beurteilen. \newline
            Medienpsychologische Theorien anwenden.
        \item Analyse\\
            Gestaltungsparameter untersuchen und Produktionsprozesse darstellen.
            \newline
            Medienpsychologische Prozesse analysieren.
    \end{itemize}
    
\newpage
    \section{Mediengestaltung}

\newpage
    \section{Medienpsychologie}

\newpage
\newpage
\chapter{Programmieren}

Dieser Kurs ist nicht für komplette Progameiranfänger ohne großen Zeitaufwand zu bewältigen. Der Eisenbiegler ist eingeborener Informatiker, jedoch leider nicht ein guter Professor. Ich empfehle euch Java unabhängig von seinen Resourcen zu lernen, seine Übungsaufgaben sind jedoch Sinnvoll konzipiert. Es steigen regelmäßig über mehrere Semester mehrere Personen, meist ohne Vorkentnisse, an der Bildverarbeitungsaufgabe aus.\\
% nicht sehr motivierend bro
Wer einen Onlinekurs nebenbei belegen will empfehle ich Kurse von \href{https://open.hpi.de/courses}{openhpi}, dies sind kostenfreie Kurse. Speziell empfehle ich diesen \href{https://open.hpi.de/courses/javaeinstieg2020}{Java Kurs}, als alternative zu diesem gibt es auch noch die  \href{https://open.hpi.de/courses/javaeinstieg-schule2024}{Schülerversion}.

\newpage
\section{SPO}

    Nachdem Studierende das Modul erfolgreich abgeschlossen haben, können sie:
    \begin{itemize}
        \item Wissen / Kenntnisse\\
            Die Sprachelemente einer imperativen Programmiersprache benennen.
        \item Verstehen\\
            Die Bedeutung eines imperativen Computerprogramms erklären.
        \item Anwenden\\
            Mit einer integrierten Entwicklungsumgebung arbeiten.
        \item Analyse\\
            Den Ablauf eines vorgegebenen imperativen Computerprogramms beschreiben.
        \item Synthesis\\
            Zu einer einfachen Aufgabenstellung ein imperatives Computerprogramm selbstständig implementieren.
        \item Evaluation\\
            Unterschiedliche Computerprogramme in Bezug auf ihre Effizienz miteinander vergleichen.
    \end{itemize}
    
\newpage
\subsection{Theorie}\newpage
\subsection{Übungsaufgaben}

\input{Programmieren/Uebungen_Eisenbiegler/A1P_Einführung}
\input{Programmieren/Uebungen_Eisenbiegler/A2P_ImperativeProg}
\input{Programmieren/Uebungen_Eisenbiegler/A3P_ImperativeProg}
\input{Programmieren/Uebungen_Eisenbiegler/A4P_Bildverarbeitung}
\input{Programmieren/Uebungen_Eisenbiegler/A5P_ImperativeProg}
\input{Programmieren/Uebungen_Eisenbiegler/A6P_Rekursion}
\input{Programmieren/Uebungen_Eisenbiegler/A7P_Objekte}
\input{Programmieren/Uebungen_Eisenbiegler/A8P_VerketteteObjekte}
\input{Programmieren/Uebungen_Eisenbiegler/A9P_ExceptionHandle}
\input{Programmieren/Uebungen_Eisenbiegler/A10P_ExceptionsAbst}
\input{Programmieren/Uebungen_Eisenbiegler/A11P_PersonenImKrankenhaus}
\input{Programmieren/Uebungen_Eisenbiegler/A12P_Fensterbestellung}
\chapter{MINT}
    \section{SPO}

    Nachdem Studierende das Modul erfolgreich abgeschlossen haben, können sie:
    \begin{itemize}
        \item Wissen / Kenntnisse\\
            Geometrische und algebraische Fragestellungen präzise mithilfe der adäquaten Fachbegriffe artikulieren. 
            \newline
            Zentrale Grundbegriffe der Optik sicher wiedergeben.
        \item Verstehen\\
            Mathematische Sinnzusammenhänge und Beweiselemente bzw. Herleitungen erkennen verstehen und wiedergeben. 
            \newline
            Mathematische Modelle physikalischer Phänomene (z.B. geometrisch-optisches paraxiales Arbeitsmodell für abbildende Systeme) verstehen.
        \item Anwenden\\
            Techniken der Vektorrechnung und der Matrixalgebra auf geometrische Probleme anwenden. \newline
            Grundgesetze der Strahlenoptik auf einfache Kameraobjektivmodelle bzw. Fragestellungen der Fotografie anwenden.
        \item Analyse\\
            Geometrische Standardprobleme in der Ebene und im Raum analysieren. 
            \newline
            Angemessen ausgewählte physikalische Systeme und Strukturen selbstständig analysieren und beschreiben.
        \newpage
        \item Synthesis\\
            Für Frage- und Problemstellungen aus (Linearer) Algebra und Geometrie unter den bereitgestellten Hilfsmitteln die jeweils adäquaten auswählen. 
            \newline
            Ein geeignetes eingegrenztes, für die Medientechnik relevantes Thema aus Optik oder Akustik im Überblick darstellen.
        \item Evaluation\\
            Verschiedene Verfahren (z.B. zur Bestimmung affiner Transformationen) hinsichtlich Übersichtlichkeit und Aufwand abwägen.
    \end{itemize}
    
\newpage
    \section{Mathe}
    \begin{itemize}
        \item {Grundlegende Kenntnisse}
        \item{Einheitskreis}
    \end{itemize}
    
    \subsection{Einleitung}
    \input{MINT/Mathe_Lasowski/Grundlegendes}
    \input{MINT/Mathe_Lasowski/Einstieg_und_Wiederholung}
\newpage
    \section{Physik}
    \begin{itemize}
        \item{Einleitung}
        \item{Optik abbildender Systeme}
    \end{itemize}

    \input{MINT/Physik_Garcia/Einleitung}
    \input{MINT/Physik_Garcia/Optik_abbildender_Systeme}
    Keine Ahnung :(
\newpage
    \input{MINT/2_MINT_Prüfungsaufgaben}
\newpage
\include{Buecherliste/Bücherliste}
% ======= Quellen %     hinzufügen vom Quellendokument

\addcontentsline{toc}{chapter}{Quellen}
\addcontentsline{toc}{section}{Weblinks}
\printbibliography[type=online,title={Weblinks}] % nur online-artikel

\addcontentsline{toc}{section}{Artikel}
\printbibliography[type=article,title={Articles only}]

\end{document}